% explainer-source-hash: 97967a9d45db4c4a54a04f7b9e4d97567a9e23aa92e39f3cd61168e4e70269c7
% explainer-source-commit: 0de7610f533d267cc97187ea2d75d1ee6ac1aab2
% explainer-generated: 2026-07-16
% Regenerate with /explain-all (see WIRING.md). Do not edit this stamp by hand:
% it is how the toolchain knows whether this document still matches the code.
\documentclass[11pt,a4paper]{article}
\input{preamble}

\title{\textbf{Flask Online Shop}\\[0.3em]\large A Deep Dive into a Flask + Stripe Storefront}
\author{Portfolio explainer, written for a reader with no prior background}
\date{}

\begin{document}
\maketitle
\thispagestyle{empty}

\begin{keyidea}{The one-paragraph version}
\textbf{Flask Online Shop} is a small web store that sells planets as a joke. A visitor
loads a page of nine product cards, can register an account, log in, and click
\emph{Buy}, which hands them off to Stripe's own payment page. It is roughly 360 lines
of Python in one file, five HTML templates, and one stylesheet. It exists to wire up the
parts a real store needs (accounts, password hashing, a database, a payment provider)
without a framework that hides how any of it works. It succeeds at that, and it also
contains a login system that is broken in a way worth understanding precisely, because
understanding exactly \emph{how} it is broken teaches more than the working parts do.
\end{keyidea}

\tableofcontents
\newpage

\section{Before anything else: the vocabulary}

This section defines every term the rest of the document leans on. Skip it if these are
already familiar.

\begin{jargon}{Web application, server, client}
When you open a web page, your browser (the \textbf{client}) sends a message across the
network asking for a page. A program on another machine (the \textbf{server}) receives
that message, decides what to say back, and sends a reply. A \textbf{web application} is
that server-side program. This project is one.
\end{jargon}

\begin{jargon}{HTTP, request, response, GET, POST}
\textbf{HTTP} is the language browsers and servers speak. The browser sends a
\textbf{request}; the server sends back a \textbf{response}. Each request carries a
\textbf{method} announcing intent. \code{GET} means "give me this page" and is what a
browser does when you type an address. \code{POST} means "here is some data, do something
with it" and is what a browser does when you submit a form. The distinction matters here:
logging in is a \code{POST} because it submits a password.
\end{jargon}

\begin{jargon}{Status code, redirect}
Every response carries a number. \code{200} means success. \code{302} and \code{303} mean
\textbf{redirect}: "do not display anything, go ask for this other address instead." The
browser follows it automatically. \code{500} means the server code crashed. This project
uses redirects after login, after registration, and to hand the visitor to Stripe.
\end{jargon}

\begin{jargon}{Web framework, Flask}
A \textbf{web framework} is a library that handles the tedious parts of being a server
(parsing HTTP, matching addresses to code, building responses) so you write only your own
logic. \textbf{Flask} is a deliberately small Python framework. It gives you routing and
templates and little else, which is exactly why it was chosen here: nothing is hidden.
\end{jargon}

\begin{jargon}{Route, endpoint, view function}
A \textbf{route} is a rule connecting a URL path to a Python function. In Flask you write
\code{@app.route('/login')} above a function, and Flask calls that function whenever a
request arrives for \code{/login}. The function is called a \textbf{view function}. This
app has seven routes.
\end{jargon}

\begin{jargon}{Template, Jinja}
A \textbf{template} is an HTML file with holes in it. \textbf{Jinja} is the engine that
fills the holes with real values. \code{\{\{ planet.name \}\}} is a hole; Jinja replaces
it with \code{Mercury}. \code{\{\% for ... \%\}} repeats a block of HTML once per item,
which is how nine product cards come from one card written once.
\end{jargon}

\begin{jargon}{Database, SQL, SQLite, table, row, column}
A \textbf{database} stores data so it survives the program shutting down. \textbf{SQL} is
the language for querying it. \textbf{SQLite} is a database that is just a single file on
disk, with no separate server to install, which is why small projects use it. Data lives
in \textbf{tables}: a table is a grid, each \textbf{row} is one record (one user), each
\textbf{column} is one field (the email).
\end{jargon}

\begin{jargon}{ORM, SQLAlchemy}
An \textbf{ORM} (Object Relational Mapper) lets you treat database rows as ordinary Python
objects instead of writing SQL by hand. \textbf{SQLAlchemy} is the standard Python ORM.
Writing \code{User.query.filter\_by(email=email).first()} produces a \code{SELECT}
statement behind the scenes and hands back a Python object.
\end{jargon}

\begin{jargon}{Hashing, bcrypt, salt}
Storing passwords as plain text is indefensible: anyone who reads the database learns
every password. A \textbf{hash} is a one-way scramble. You can compute the hash of a
password, but you cannot reverse the hash back into the password. To check a login you
hash what was typed and compare hashes. \textbf{bcrypt} is a hashing function designed to
be deliberately \emph{slow}, so an attacker guessing billions of passwords is throttled. A
\textbf{salt} is random data mixed in, so two users with the same password still get
different hashes. bcrypt generates and embeds the salt for you.
\end{jargon}

\begin{jargon}{Session, cookie}
HTTP has no memory: each request arrives knowing nothing about the last. A \textbf{cookie}
is a small piece of data the server hands the browser, which the browser sends back on
every later request. That is how a site remembers who you are. A \textbf{session} is the
server-side notion of "this particular visitor, across their requests," normally carried
by a cookie. \textbf{Remember this definition.} The central defect in this project is that
it tracks login \emph{without} using the session at all.
\end{jargon}

\begin{jargon}{Environment variable}
A value handed to a program by the operating system when it starts, rather than written
into the source code. Secrets belong here, because source code gets committed to version
control and shared, while environment variables do not.
\end{jargon}

\begin{jargon}{Stripe, Stripe Checkout}
\textbf{Stripe} is a payment company. \textbf{Stripe Checkout} is its hosted payment page:
instead of building a card form (and inheriting the legal and security burden of touching
card numbers), the app asks Stripe to create a \textbf{Checkout Session}, then redirects
the buyer to a page on Stripe's own domain. The card number never touches this
application. That is the single best security decision in the project.
\end{jargon}

\newpage
\section{Orientation: what the project physically is}

\subsection{The files}

Every tracked file, with its role:

\begin{figure}[H]
\centering
\begin{forest} dirtree
[flask-online-shop
  [main.py \quad{\rmfamily\itshape the entire application: config, model, catalog, all 7 routes}]
  [requirements.txt \quad{\rmfamily\itshape 5 pinned dependencies}]
  [.env.example \quad{\rmfamily\itshape documents the 3 environment variables}]
  [.gitignore \quad{\rmfamily\itshape excludes .env, instance/, *.sqlite3}]
  [README.md]
  [LICENSE \quad{\rmfamily\itshape PolyForm Strict: viewable, not reusable}]
  [templates/ \quad{\rmfamily\itshape Jinja HTML}
    [base.html \quad{\rmfamily\itshape shared shell: nav, flash messages, footer}]
    [index.html \quad{\rmfamily\itshape the catalog, one card looped 9 times}]
    [login.html]
    [create\_user.html]
    [success.html]
    [cancel.html]
  ]
  [static/
    [css/styles.css \quad{\rmfamily\itshape 321 lines, CSS variables + card components}]
    [images/ \quad{\rmfamily\itshape mercury.jpeg ... pluto.jpeg, 9 files}]
  ]
  [docs/
    [screenshots/app.png]
    [explainers/ \quad{\rmfamily\itshape this document}]
  ]
]
\end{forest}
\caption{The complete repository. There is no test suite, no CI configuration, and no
migration folder; every one of those absences is discussed later.}
\end{figure}

\begin{keyidea}{Everything is in one file, and that is a real decision}
\file{main.py} holds the configuration, the database model, the product catalog, and all
seven routes. For a project this size that is defensible and even helpful: a reader sees
the whole system in one scroll. It is the first thing that would have to change if the
project grew, and the natural split (config, models, routes, catalog data) is obvious.
\end{keyidea}

\subsection{The dependencies, and what each one is actually for}

\begin{table}[H]
\centering\small
\begin{tabular}{@{}llp{7.7cm}@{}}
\toprule
\textbf{Package} & \textbf{Pin} & \textbf{Why it is here} \\
\midrule
Flask & 3.0.3 & The web framework itself: routing, templates, requests, redirects. \\
Flask-SQLAlchemy & 3.1.1 & Glue that wires SQLAlchemy into Flask's app lifecycle. \\
SQLAlchemy & 2.0.30 & The ORM. Turns the \code{User} class into a SQLite table. \\
stripe & 9.12.0 & Official Stripe SDK. Used for exactly one call, to create a Checkout Session. \\
bcrypt & 4.1.3 & Password hashing and verification. \\
\bottomrule
\end{tabular}
\caption{All five dependencies are pinned to exact versions, which is good practice: an
unpinned install can silently pull a breaking release months later.}
\end{table}

\begin{honest}{The dependency list is missing something the README assumes}
The README instructs the reader to \code{cp .env.example .env} and edit it. Nothing in
this project ever reads a \file{.env} file. There is no \code{python-dotenv} in
\file{requirements.txt} and no \code{load\_dotenv()} call anywhere in the source; I
searched the whole tree for both and found neither. Flask does not read \file{.env} on its
own (that is a feature of the \code{flask} CLI with \code{python-dotenv} installed, and
this app is launched as \code{python main.py}). So a reader who follows the setup
instructions exactly will create a \file{.env} file that is silently ignored, get the
insecure fallback signing key and an empty Stripe key, and only discover it when checkout
fails. The variables must instead be \code{export}ed into the shell. This is a genuine
README-versus-code disagreement, and the code is the truth.
\end{honest}

\newpage
\section{Architecture}

\subsection{The whole system in one picture}

\begin{figure}[H]
\centering
\resizebox{\textwidth}{!}{
\begin{tikzpicture}[node distance=1cm]
  \node[box, minimum width=2.6cm] (browser) {Browser\\{\scriptsize the visitor}};
  \node[box, right=3.2cm of browser, minimum width=3.6cm, minimum height=3.4cm] (flask)
    {\textbf{Flask app}\\{\scriptsize main.py}\\[2pt]{\scriptsize 7 routes}\\{\scriptsize context processor}};
  \node[store, right=3.2cm of flask] (db) {SQLite\\{\scriptsize instance/}\\{\scriptsize db.sqlite3}};
  \node[extbox, above=1.6cm of db, minimum width=2.8cm] (stripe) {Stripe\\{\scriptsize hosted Checkout}};
  \node[extbox, below=1.6cm of flask, minimum width=3.2cm] (tpl) {templates/ + static/\\{\scriptsize Jinja, CSS, 9 images}};

  \draw[flow] ([yshift=6mm]browser.east) -- node[lbl, above]{HTTP request} ([yshift=6mm]flask.west);
  \draw[flow] ([yshift=-6mm]flask.west) -- node[lbl, below]{HTML response} ([yshift=-6mm]browser.east);
  \draw[flow] (flask.east) -- node[lbl, above, align=center]{ORM queries\\via Flask-SQLAlchemy} (db.west);
  \draw[flow] (flask.north east) -- node[lbl, above left, align=center]{create\\Checkout Session} (stripe.west);
  \draw[dflow] (browser.north) |- node[lbl, above, pos=0.75, align=center]{redirect 303: buyer enters card\\ON STRIPE'S DOMAIN, never here} (stripe.north west);
  \draw[flow] (flask.south) -- node[lbl, right]{render} (tpl.north);
\end{tikzpicture}}
\caption{The system. Solid arrows are this application's own traffic; the dashed arrow is
the buyer's browser leaving for Stripe. Card details flow only along the dashed path, so
they never enter this codebase, this server, or this database.}
\end{figure}

\begin{keyidea}{Why the dashed arrow is the most important line in the diagram}
Handling raw card numbers puts an application under the full weight of PCI DSS (the card
industry's security standard) and makes any breach catastrophic. By redirecting to a page
Stripe hosts, this project never sees a card number. A beginner project that built its own
card form would be far more impressive-looking and far more dangerous. This is the right
call, and it is worth being able to say why in an interview.
\end{keyidea}

\subsection{The seven routes}

\begin{table}[H]
\centering\small
\begin{tabular}{@{}p{3.5cm}lp{3.6cm}p{5.2cm}@{}}
\toprule
\textbf{Path} & \textbf{Methods} & \textbf{Function} & \textbf{What it does} \\
\midrule
\code{/} & GET & \code{home} & Renders the catalog from \code{get\_planets()}. \\
\code{/create-checkout-\allowbreak session} & POST & \code{create\_checkout\_\allowbreak session} & Builds a Stripe session for one planet, redirects to Stripe. \\
\code{/success} & GET & \code{success} & Static "payment worked" page. \\
\code{/cancel} & GET & \code{cancel} & Static "cancelled" page. \\
\code{/login} & GET, POST & \code{login} & Shows the form; checks the password. \\
\code{/logout} & GET & \code{logout} & Clears a logged-in flag. \\
\code{/add\_user} & GET, POST & \code{add\_user} & Shows the form; creates an account. \\
\bottomrule
\end{tabular}
\caption{Note the inconsistent naming: \code{/create-checkout-session} uses hyphens while
\code{/add\_user} uses an underscore. Cosmetic, but a reviewer will notice.}
\end{table}

\newpage
\section{The catalog: how a product becomes a card}

\subsection{Two dictionaries and the function that joins them}

Product data lives in two Python dictionaries at the top of \file{main.py}:

\begin{itemize}
  \item \code{planet\_price\_dict} maps a slug (\code{'mercury'}) to a structure shaped
  exactly like a \textbf{Stripe line item}: currency, product name, and
  \code{unit\_amount}, which Stripe requires in \emph{cents}.
  \item \code{PLANET\_INFO} maps the same slug to display copy: name, image filename, and
  a description.
\end{itemize}

\code{get\_planets()} walks \code{PLANET\_ORDER} and joins the two into the list the
template loops over.

\begin{figure}[H]
\centering
\resizebox{\textwidth}{!}{
\begin{tikzpicture}[node distance=0.9cm]
  \node[box, minimum width=3.3cm] (ppd) {\code{planet\_price\_dict}\\{\scriptsize slug $\rightarrow$ line item}\\{\scriptsize unit\_amount in CENTS}};
  \node[box, below=1.6cm of ppd, minimum width=3.3cm] (pi) {\code{PLANET\_INFO}\\{\scriptsize slug $\rightarrow$ name,}\\{\scriptsize image, copy}};
  \node[box, right=1.7cm of ppd, minimum width=2.7cm] (fp) {\code{format\_price()}\\{\scriptsize cents $\rightarrow$ "\$42,069"}};
  \node[box, right=1.7cm of fp, minimum width=2.5cm, minimum height=2.0cm] (gp) {\code{get\_planets()}\\{\scriptsize joins both}};
  \node[box, right=1.7cm of gp, minimum width=2.5cm] (tpl) {\file{index.html}\\{\scriptsize Jinja for-loop}};
  \node[extbox, above=1.4cm of gp, minimum width=2.7cm] (stripe) {Stripe API\\{\scriptsize same number}};

  \draw[flow] (ppd.east) -- node[lbl, above]{cents} (fp.west);
  \draw[flow] (fp.east) -- node[lbl, above, align=center]{price\_\\display} (gp.west);
  \draw[flow] (pi.east) -| node[lbl, below, pos=0.25, align=center]{name, image, description} (gp.south);
  \draw[flow] (gp.east) -- node[lbl, above, align=center]{list of 9\\dicts} (tpl.west);
  \draw[flow] (ppd.north) |- node[lbl, above, pos=0.65, align=center]{the SAME dict becomes the line item} (stripe.west);
\end{tikzpicture}}
\caption{One number, two destinations. The price shown on the card and the price charged
by Stripe are derived from the same \code{unit\_amount}, so they cannot drift apart.}
\end{figure}

\begin{keyidea}{Why this indirection is not over-engineering}
The README records that the template once hardcoded "Uranus: \$69,420" while the dict
actually sent to Stripe said \code{4206900} cents, which is \$42,069. A customer would
have seen one number and been charged another. That is a bug class, not a typo: the same
fact was typed in two places. Deriving the display from the charge makes the bug
unrepresentable. I verified the fix holds: the running app renders \$42,069 for Uranus,
matching the Stripe amount exactly.
\end{keyidea}

\subsection{\code{format\_price()}, precisely}

\begin{lstlisting}[language=Python, caption={The whole function. Cents in, display string out.}]
def format_price(unit_amount_cents):
    dollars = unit_amount_cents / 100
    if dollars == int(dollars):
        return f'${int(dollars):,}'
    return f'${dollars:,.2f}'
\end{lstlisting}

The \code{:,} in the f-string inserts thousands separators. The branch drops trailing
cents on whole-dollar amounts, so Mercury reads \$10,000 rather than \$10,000.00, while
Earth keeps its cents at \$777,777.77. I ran \code{get\_planets()} directly and confirmed
all nine outputs; they match the underlying cent amounts in every case.

\begin{honest}{Pluto almost certainly cannot be bought}
Pluto's \code{unit\_amount} is \code{999999999999999} cents, which the page renders as
\$9,999,999,999,999.99. Stripe's documented maximum \code{unit\_amount} for USD is
\code{99999999} cents (\$999,999.99). Pluto's amount exceeds that ceiling by roughly ten
million times, so the Stripe API would be expected to reject the Checkout Session with an
invalid-amount error, meaning the button on "the people's favorite" fails while the other
eight work.

\textbf{This is labelled inferred, not verified.} Confirming it requires a real Stripe
test key, which this environment does not have, so I could not observe the actual API
response. The claim rests on Stripe's published limit and on Pluto's amount, which I did
read directly from the source. It is a joke value that was never exercised, and it is
exactly the kind of thing an interviewer might click first. Earth (\$777,777.77) sits
safely under the ceiling; only Pluto exceeds it.
\end{honest}

\subsection{The template loop}

\begin{lstlisting}[language=HTML, caption={\file{index.html}: nine cards from one block of markup.}]
{% for planet in planets %}
    <article class="product-card">
        <img class="product-image" src="{{ url_for('static', filename='images/' + planet.image) }}" alt="{{ planet.name }}">
        <div class="product-body">
            <h3 class="product-name">{{ planet.name }}</h3>
            <div class="product-price">{{ planet.price_display }}</div>
            <p class="product-description">{{ planet.description }}</p>
            <form action="{{ url_for('create_checkout_session') }}" method="POST">
                <button name="name" value="{{ planet.slug }}" type="submit" class="btn btn-buy">Buy {{ planet.name }}</button>
            </form>
        </div>
    </article>
{% endfor %}
\end{lstlisting}

The README records that this replaced nine copy-pasted card blocks. The mechanism of the
Buy button is worth noting: the \code{<button>} carries \code{name="name"} and
\code{value="\{\{ planet.slug \}\}"}, so submitting the form POSTs a single field called
\code{name} whose value is the slug. That is what the checkout route reads.

\newpage
\section{The data model}

\subsection{One table, four columns}

\begin{lstlisting}[language=Python, caption={The entire persistence layer.}]
class User(db.Model):
    id = db.Column(db.Integer, primary_key=True)
    email = db.Column(db.String(100))
    password = db.Column(db.String(100))
    logged_in = db.Column(db.Integer)
\end{lstlisting}

\begin{figure}[H]
\centering
\begin{tikzpicture}
  \node[box, minimum width=11cm, minimum height=1cm] (t) {\textbf{user} \quad{\scriptsize (SQLite table, auto-created at startup)}};
  \node[box, below=0.35cm of t.south west, anchor=north west, minimum width=2.4cm] (a) {\code{id}\\{\scriptsize INTEGER, PK}};
  \node[box, right=0.2cm of a, minimum width=2.4cm] (b) {\code{email}\\{\scriptsize String(100)}};
  \node[box, right=0.2cm of b, minimum width=2.9cm] (c) {\code{password}\\{\scriptsize String(100)}};
  \node[box, right=0.2cm of c, minimum width=2.6cm] (d) {\code{logged\_in}\\{\scriptsize Integer, 0 or 1}};
\end{tikzpicture}
\caption{The complete schema. There are no other tables: no products (they are hardcoded
in Python), no orders, no payments. Nothing about a purchase is ever recorded.}
\end{figure}

\begin{jargon}{Primary key}
A column whose value uniquely identifies a row. \code{primary\_key=True} on \code{id} makes
SQLite assign each new user an automatically increasing number.
\end{jargon}

\begin{honest}{The schema has four separate defects}
\begin{enumerate}
  \item \textbf{\code{password} is declared \code{String(100)} but stores bytes.}
  \code{bcrypt.hashpw()} returns \code{bytes}, and that is what gets assigned. I queried
  the live database with \code{typeof(password)} and SQLite reports \code{blob}, not
  \code{text}. SQLite tolerates this because it is dynamically typed; a stricter database
  such as PostgreSQL would reject or mangle it. The consequence is visible in the login
  route, which carries an \code{isinstance(user.password, str)} guard to handle a value
  that may come back as either type. That guard is a symptom of the wrong column type, not
  a fix for it. The column should be \code{LargeBinary}, or the hash should be decoded to
  a string before storage.
  \item \textbf{\code{email} has no uniqueness constraint.} Nothing stops two rows sharing
  an address. I registered \code{dup@x.com} twice against the running app and both
  inserts succeeded, giving two rows with the same email. Since login does
  \code{filter\_by(email=...).first()}, the second account is permanently unreachable: its
  owner can never log in, because \code{first()} always returns row one. The fix is one
  argument: \code{unique=True}.
  \item \textbf{\code{logged\_in} is the wrong place for that fact entirely.} It is
  discussed at length in the next section; it is the project's central defect.
  \item \textbf{No length or format validation.} \code{String(100)} is not enforced by
  SQLite, and nothing checks the email is an email.
\end{enumerate}
\end{honest}

\subsection{Where the database file lives, and why that changed}

\begin{lstlisting}[language=Python]
base_path = os.path.dirname(os.path.abspath(__file__))
instance_dir = os.path.join(base_path, 'instance')
os.makedirs(instance_dir, exist_ok=True)
path_to_sql_file = os.path.join(instance_dir, 'db.sqlite3')
\end{lstlisting}

The README records that this path was once built from \code{sys.argv[0]}, which resolves
to whatever launched the process and therefore moved depending on how the app was started.
Deriving it from \code{\_\_file\_\_} (the location of the source file itself) is stable
regardless of the working directory. \code{db.create\_all()} then runs inside an
\textbf{application context} at import time and creates the table if it is absent; it is
idempotent, so calling it on every startup is safe.

\begin{jargon}{Application context}
Flask keeps certain globals (the current app, the database handle) in a context that must
be explicitly active outside a request. \code{with app.app\_context():} is what makes
\code{db.create\_all()} legal at module level.
\end{jargon}

\begin{honest}{\code{create\_all()} is not a migration system}
It creates tables that do not exist. It will never alter one that does. Add
\code{unique=True} to \code{email} tomorrow and existing databases keep the old schema
silently, with no error and no upgrade. Real projects use a migration tool such as Alembic.
For a demo this is an acceptable trade, but it is a limit worth naming out loud rather
than being caught by.
\end{honest}

\newpage
\section{Authentication, and the defect at the centre of it}

\subsection{Registration: \code{/add\_user}}

\begin{figure}[H]
\centering
\begin{tikzpicture}[node distance=0.8cm]
  \node[box, minimum width=3.1cm] (form) {POST \code{/add\_user}\\{\scriptsize email, password,}\\{\scriptsize re-entered-pw}};
  \node[dec, right=1.9cm of form, text width=1.7cm] (match) {pw == re-pw?};
  \node[box, right=1.9cm of match, minimum width=3.0cm] (salt) {\code{bcrypt.gensalt()}\\{\scriptsize random salt}};
  \node[box, below=1.5cm of match, minimum width=3.4cm] (flash) {flash "passwords didn't match"\\{\scriptsize redirect back to form}};
  \node[box, below=1.5cm of salt, minimum width=3.0cm] (hash) {\code{bcrypt.hashpw()}\\{\scriptsize slow by design}};
  \node[store, below=1.4cm of hash] (db) {user row\\{\scriptsize insert}};
  \node[box, left=2.4cm of db, minimum width=2.6cm] (home) {redirect \code{/}};

  \draw[flow] (form) -- (match);
  \draw[flow] (match) -- node[lbl, right]{no} (flash);
  \draw[flow] (match) -- node[lbl, above]{yes} (salt);
  \draw[flow] (salt) -- node[lbl, right]{salt} (hash);
  \draw[flow] (hash) -- node[lbl, right, align=center]{hash bytes} (db);
  \draw[flow] (db) -- (home);
\end{tikzpicture}
\caption{The registration flow. Note what is absent: no check that the email is already
taken, no minimum password length, no email format validation, and no login afterwards.}
\end{figure}

\begin{keyidea}{The password is never stored, and that is the point}
Only \code{bcrypt.hashpw(pw\_to\_bytes, bcrypt.gensalt())} reaches the database. If someone
steals \file{db.sqlite3}, they get hashes. Because bcrypt is intentionally slow and each
hash carries its own random salt, cracking them is expensive and cannot be done in bulk
with a precomputed table. This part of the project is done correctly.
\end{keyidea}

\begin{honest}{Registration validation is thinner than it looks}
The registration form declares \code{maxlength="10"} on both password fields. That is an
HTML attribute, enforced only by the browser. I posted a 60-character password directly to
\code{/add\_user}, bypassing the form: it was accepted, stored, and I then logged in with
it successfully. The README lists this honestly under "What I'd do differently," and the
observation is correct. Two related crashes I confirmed against the running app:

\begin{itemize}
  \item Posting \code{/add\_user} with the \code{password} field \textbf{missing entirely}
  returns HTTP 500. \code{request.form.get('password')} returns \code{None} and
  \code{bytes(None, 'UTF-8')} raises \code{TypeError}. A form always sends the field, so a
  browser user never sees this, but any direct client triggers it.
  \item There is no server-side email validation whatsoever; \code{type="email"} is,
  again, browser-only.
\end{itemize}
A rule worth internalising: \textbf{every client-side validation is a suggestion.} The
server is the only thing that can enforce anything.
\end{honest}

\subsection{Login: \code{/login}}

\begin{lstlisting}[language=Python, caption={The heart of the login route.}]
user = User.query.filter_by(email=email).first()
if user:
    encoded_pw = bytes(pw, 'UTF-8')
    check = bcrypt.checkpw(encoded_pw, user.password.encode('UTF-8') if isinstance(user.password, str) else user.password)
    if check:
        flash('You are now Logged In')
        user.logged_in = 1
        db.session.commit()
        return redirect(url_for('home'))
    else:
        error = 'That password is incorrect'
        return render_template('login.html', error=error)
else:
    error = 'That email is not in our system'
    return render_template('login.html', error=error)
\end{lstlisting}

The README describes a previous version that looked the submitted email up, then always
compared against user id 1's hash regardless of who was logging in. With one user that
works by coincidence; with two it lets the wrong password through or rejects the right
one. The current code looks up the row \emph{by the submitted email} and checks against
\emph{that row's own hash}, which is correct.

I verified this behaves as claimed by driving the running server with two real accounts:

\begin{table}[H]
\centering\small
\begin{tabular}{@{}lll@{}}
\toprule
\textbf{Attempt} & \textbf{Expected} & \textbf{Observed} \\
\midrule
Alice logs in with Alice's password & accepted & accepted, redirect to \code{/} \\
Bob logs in with Bob's password & accepted & accepted, redirect to \code{/} \\
Bob logs in with Alice's password & rejected & "That password is incorrect" \\
Unknown email & rejected & "That email is not in our system" \\
\bottomrule
\end{tabular}
\caption{The per-user hash check is genuinely fixed. Verified against a live server.}
\end{table}

\begin{honest}{Two smaller issues in this route}
\textbf{Username enumeration.} The route distinguishes "that email is not in our system"
from "that password is incorrect." That tells an attacker which addresses have accounts,
which is useful for targeting. Production systems return one generic message for both.
For a joke storefront the stakes are nil, but the reasoning is worth knowing.

\textbf{The \code{isinstance} guard.} The line is long and does two jobs at once. It exists
purely because the column type is wrong (see the schema section). Fixing the column would
delete the guard.

\textbf{No rate limiting.} Nothing throttles login attempts. bcrypt's slowness is the only
brake on guessing.
\end{honest}

\subsection{The central defect: login is a global flag, not a session}

This is the most important thing in the codebase. Two functions carry it:

\begin{lstlisting}[language=Python, caption={Six lines that decide who is logged in.}]
def is_logged_in():
    return bool(db.session.query(exists().where(User.logged_in == 1)).scalar())

@app.context_processor
def inject_nav_state():
    return {'logged_in': is_logged_in()}
\end{lstlisting}

Read \code{is\_logged\_in()} carefully. It asks the database a single question: \textbf{does
there exist any row, anywhere, whose \code{logged\_in} column equals 1?} It does not
receive, and cannot receive, any information about \emph{who is asking}. There is no
cookie read, no session lookup, no user identity involved at all.

\begin{jargon}{Context processor}
A Flask hook that injects variables into every template render automatically. Here it makes
\code{logged\_in} available to \file{base.html} on every page, so the nav link is
consistent site-wide. The mechanism is used correctly; the value it injects is the problem.
\end{jargon}

\begin{figure}[H]
\centering
\resizebox{\textwidth}{!}{
\begin{tikzpicture}[node distance=1cm]
  \node[box, minimum width=2.7cm] (alice) {Alice\\{\scriptsize logs in properly}};
  \node[box, below=0.7cm of alice, minimum width=2.7cm] (bob) {Bob\\{\scriptsize never logged in}};
  \node[box, below=0.7cm of bob, minimum width=2.7cm] (stranger) {Total stranger\\{\scriptsize no cookies at all}};
  \node[box, right=3.4cm of bob, minimum width=3.6cm, minimum height=2.6cm] (fn)
    {\code{is\_logged\_in()}\\[3pt]{\scriptsize "does ANY row}\\{\scriptsize have logged\_in == 1?"}\\[3pt]{\scriptsize identity is never consulted}};
  \node[store, right=2.8cm of fn] (db) {user table\\{\scriptsize Alice: 1}\\{\scriptsize Bob: 0}};
  \node[box, below=1.5cm of db, minimum width=3.4cm] (out) {Every visitor sees\\\textbf{"Logout"}};

  \draw[flow] (alice.east) -- node[lbl, above, pos=0.6]{GET /} (fn.west);
  \draw[flow] (bob.east) -- (fn.west);
  \draw[flow] (stranger.east) -- node[lbl, below, pos=0.6]{GET /} (fn.west);
  \draw[flow] (fn.east) -- node[lbl, above]{SELECT EXISTS} (db.west);
  \draw[flow] (db.south) -- node[lbl, right, align=center]{true} (out.north);
\end{tikzpicture}}
\caption{The defect, drawn. Three different visitors, one shared answer. Because the
question never mentions who is asking, the answer cannot differ between them.}
\end{figure}

\begin{honest}{Verified experimentally: one user logging in logs everyone in}
I ran the application and drove it with three independent HTTP clients, each with its own
cookie jar. Results:

\begin{table}[H]
\centering\small
\begin{tabular}{@{}lll@{}}
\toprule
\textbf{Visitor} & \textbf{Before Alice logs in} & \textbf{After Alice logs in} \\
\midrule
Alice's client & shows "Login" & shows "Logout" \\
Bob's client (never logged in) & shows "Login" & \textbf{shows "Logout"} \\
Brand-new client, zero cookies & shows "Login" & \textbf{shows "Logout"} \\
\bottomrule
\end{tabular}
\end{table}

A visitor who has never authenticated, holding no cookie, is presented as logged in the
moment any other person in the world logs in. The README describes this as "no real
session isolation between browser tabs or devices," which \textbf{understates it}. It is
not a tab or device issue: it is that login state is a property of the \emph{application},
not of a visitor. Tabs and devices are the symptom; the shared global flag is the disease.

Two further consequences I confirmed:
\begin{itemize}
  \item \textbf{\code{/logout} logs out a stranger.} The route calls
  \code{User.query.\allowbreak filter\_by(\allowbreak logged\_in=1)\allowbreak .first()}
  and clears that row, so it logs out whichever matching row SQLite happens to return
  first, which is not necessarily the person who clicked. With Alice and Bob both flagged, one anonymous request to \code{/logout} cleared
  \emph{Alice's} row. \code{/logout} also requires no authentication of any kind.
  \item \textbf{The flag never resets.} It lives in the database, so it survives restarts.
  If the process dies while a user is flagged, the app comes back up believing someone is
  logged in, forever, until someone hits \code{/logout}.
\end{itemize}
\end{honest}

\begin{keyidea}{Why this happened, and the honest silver lining}
Flask's \code{session} object is imported at the top of \file{main.py} and is used exactly
once in the whole file, at \code{session.pop('\_flashes', None)}, to clear stale flash
messages. The correct fix is to put the user's id in that session
(\code{session['user\_id'] = user.id}), because the session is backed by a signed cookie
and is therefore per-visitor by construction. The \code{logged\_in} column would then be
deleted outright. The \code{flask-login} extension does this and more, which is what the
README proposes.

The silver lining is real: the surrounding architecture is right. The context processor is
the correct mechanism for site-wide nav state, and it takes its value from a single
function. Change those six lines to read the session and the defect disappears without
touching a template. The bug is severe but it is also \emph{localised}, which is a
consequence of the code being reasonably structured.
\end{keyidea}

\newpage
\section{Payments}

\subsection{The checkout route}

\begin{lstlisting}[language=Python, caption={The entire payment integration.}]
@app.route('/create-checkout-session', methods=['POST'])
def create_checkout_session():
    planet = request.form.get('name')
    checkout_session = stripe.checkout.Session.create(
        line_items=[planet_price_dict[planet]],
        mode='payment',
        success_url=f'{YOUR_DOMAIN}/success',
        cancel_url=f'{YOUR_DOMAIN}/cancel',
    )
    return redirect(checkout_session.url, code=303)
\end{lstlisting}

\begin{figure}[H]
\centering
\resizebox{\textwidth}{!}{
\begin{tikzpicture}[node distance=0.9cm]
  \node[box, minimum width=2.6cm, minimum height=1.2cm] (b) {Browser};
  \node[box, right=3.4cm of b, minimum width=2.8cm, minimum height=1.2cm] (app) {Flask\\{\scriptsize checkout route}};
  \node[extbox, right=3.4cm of app, minimum width=2.8cm, minimum height=1.2cm] (api) {Stripe API\\{\scriptsize server to server}};
  \node[extbox, below=2.4cm of b, minimum width=3.2cm, minimum height=1.2cm] (pay) {Stripe hosted page\\{\scriptsize card entered HERE}};
  \node[box, below=2.4cm of app, minimum width=2.8cm, minimum height=1.2cm] (succ) {\code{/success}\\{\scriptsize static page}};

  \draw[flow] ([yshift=3mm]b.east) -- node[lbl, above, align=center]{POST\\name=earth} ([yshift=3mm]app.west);
  \draw[flow] ([yshift=-3mm]app.west) -- node[lbl, below, align=center]{303 redirect\\to Stripe} ([yshift=-3mm]b.east);
  \draw[flow] ([yshift=3mm]app.east) -- node[lbl, above, align=center]{Session.create\\line\_items} ([yshift=3mm]api.west);
  \draw[flow] ([yshift=-3mm]api.west) -- node[lbl, below, align=center]{session.url} ([yshift=-3mm]app.east);
  \draw[dflow] (b.south) -- node[lbl, left, align=center]{buyer goes to\\Stripe and pays} (pay.north);
  \draw[dflow] (pay.east) -- node[lbl, above, align=center]{redirect back\\(browser controlled)} (succ.west);
\end{tikzpicture}}
\caption{The checkout handoff. The application talks to Stripe's API server-to-server to
create the session, then sends the browser to Stripe's own page. Note the return path:
\code{/success} is reached by a plain browser redirect.}
\end{figure}

\begin{jargon}{HTTP 303}
\code{redirect(..., code=303)} is deliberate. 303 tells the browser "your POST is done,
now GET this other page." Without it a browser refresh could resubmit the POST and start a
second checkout. Using 303 after a form submission is a small, correct detail.
\end{jargon}

\subsection{Security posture of the payment flow}

\begin{honest}{The payment flow is honest about cards and dishonest about outcomes}
\textbf{What is right:} card data never touches this app (Stripe hosts the form); the
secret key comes from the environment; the amount is server-controlled, so a buyer cannot
tamper with the price by editing the form (the client sends only a slug, and the server
looks the amount up).

\textbf{What is wrong, verified against the running app:}
\begin{enumerate}
  \item \textbf{\code{/success} is a static page anyone can visit.} I requested
  \code{/success} directly, having paid nothing: HTTP 200, "payment successful." Nothing
  verifies a payment occurred. The route takes no session id and asks Stripe nothing.
  \item \textbf{There is no webhook, so the app never learns about any payment.} A
  \textbf{webhook} is a callback Stripe makes to your server, server-to-server, to confirm
  a payment truly completed. Without one, this app's only signal is the buyer's browser
  arriving at a URL, which the buyer fully controls. The README lists this correctly.
  \item \textbf{Nothing is recorded anyway.} There is no orders table. Even a real payment
  leaves no trace in this system, so "who bought what" is unanswerable. This is the deeper
  reason the missing webhook is not a small gap: there is no place to put the answer.
  \item \textbf{No login is required to buy.} The checkout route has no auth check at all.
  I posted to it while completely logged out and it proceeded to the Stripe call. The
  entire account system is therefore decorative with respect to purchasing.
  \item \textbf{An unknown planet name crashes the route.} \code{planet\_price\_dict[planet]}
  is an unguarded dict lookup. I posted \code{name=krypton} and got HTTP 500 from the
  resulting \code{KeyError}. A missing \code{name} field does the same. It should return
  400.
\end{enumerate}
\end{honest}

\begin{keyidea}{The takeaway an interviewer would want}
"Redirect to Stripe, then trust the redirect back" is the classic beginner payment
mistake, and this project makes it. The fix is a webhook plus an orders table: Stripe tells
your server directly that session \code{cs\_x} is paid, you record it, and \code{/success}
displays only what your own database confirms. Being able to explain that in one breath is
worth more than having built it.
\end{keyidea}

\newpage
\section{Configuration and secrets}

\begin{lstlisting}[language=Python]
SECRET_KEY = os.environ.get('FLASK_SECRET_KEY', 'dev-only-insecure-secret-key')
STRIPE_SECRET_KEY = os.environ.get('STRIPE_SECRET_KEY', '')
YOUR_DOMAIN = os.environ.get('YOUR_DOMAIN', 'http://127.0.0.1:5000')
\end{lstlisting}

The README records that a real Stripe test key and a Flask signing key were once hardcoded
in the source. Both now come from the environment, \file{.env} is gitignored, and
\file{.env.example} documents each variable without containing a real value. That is the
correct pattern.

\begin{jargon}{Why the signing key matters}
Flask signs the session cookie with \code{SECRET\_KEY} so a visitor cannot forge one. If
the key is publicly known, anyone can mint cookies the server will trust. Here the session
carries only flash messages, so the blast radius is small today; the moment login moves
into the session (the correct fix above), that key becomes the thing protecting every
account.
\end{jargon}

\begin{honest}{Three loose ends in configuration}
\begin{enumerate}
  \item \textbf{The insecure fallback is a footgun.} \code{'dev-only-insecure-secret-key'}
  means a deployment that forgets the variable does not crash; it runs with a publicly
  known signing key. The comment says "only for a quick local run," which is honest, but
  failing loudly in production is safer than defaulting quietly.
  \item \textbf{\file{.env} is never loaded} (detailed earlier). The setup instructions do
  not work as written.
  \item \textbf{\code{app.config['SESSION\_TYPE'] = 'filesystem'} does nothing.} It is set
  in the \code{\_\_main\_\_} block and is a \code{Flask-Session} setting. Flask-Session is
  not in \file{requirements.txt} and is never imported, so plain Flask ignores the key
  entirely. It is dead configuration, and a leftover of the same confusion that produced
  the global login flag.
\end{enumerate}
\end{honest}

\section{The front end, briefly}

\file{base.html} is the shared shell: it loads Bootstrap 5 from a CDN, then
\file{styles.css}, renders the nav (whose Login/Logout link is driven by the context
processor), renders flashed messages, and defines the \code{content} block every other
template fills. \file{styles.css} is 321 lines built on CSS variables
(\code{--color-bg}, \code{--color-primary}, and so on) with reusable \code{product-card},
\code{auth-card} and \code{result-card} components. It is tidy and consistent.

\begin{jargon}{Flash message}
Flask's mechanism for a one-shot notification that survives exactly one redirect: call
\code{flash('...')}, and the next rendered page displays it once. It is stored in the
session cookie, which is why \code{session.pop('\_flashes', None)} at the top of
\code{/login} exists: it clears any stale message before showing the form.
\end{jargon}

\begin{honest}{Small front-end notes}
\code{<script src="https://js.stripe.com/v3/"></script>} is loaded on every page but
never used: the redirect approach needs no client-side Stripe library. Bootstrap's CSS is
pinned to 5.0.2 while its JavaScript bundle is pinned to 5.3.0-alpha3, two different
versions of the same library, and neither is actually needed by anything on the page. Both
are dead weight rather than bugs.
\end{honest}

\section{Running it}

\begin{lstlisting}[language=bash, caption={What actually works, as opposed to what the README says.}]
python3 -m venv venv
source venv/bin/activate
pip install -r requirements.txt

# The README says to copy .env and edit it. Nothing reads .env.
# Export the variables instead:
export FLASK_SECRET_KEY=$(python -c "import secrets; print(secrets.token_hex(32))")
export STRIPE_SECRET_KEY=sk_test_your_key_here

python main.py     # http://127.0.0.1:5000/
\end{lstlisting}

\subsection{What I verified myself}

I installed the pinned dependencies in a clean virtual environment and drove the running
server. Confirmed: it boots; \code{GET /} returns 200 and renders exactly nine product
cards; all nine displayed prices match their Stripe cent amounts; registration, login,
wrong-password, unknown-email and logout all behave as the README claims; the per-user
hash fix is genuine. Also confirmed, and less flattering: one user logging in makes every
visitor appear logged in; duplicate emails are accepted and orphan the second account; a
60-character password bypasses the form's \code{maxlength}; \code{/success} serves 200 to
someone who never paid; an unknown planet slug and a missing password field each return
500; and checkout proceeds with no login.

\subsection{What I could not verify}

A real Stripe checkout end to end needs a test-mode secret key that this environment does
not have. With no key, \code{POST /create-checkout-session} reaches Stripe and fails on
authentication, which confirms the route is wired correctly up to that boundary and no
further. The Pluto amount claim is likewise inferred from Stripe's documented limit rather
than observed.

\section{Honest overall assessment}

\begin{table}[H]
\centering\small
\begin{tabular}{@{}p{5.4cm}p{3.1cm}p{6.3cm}@{}}
\toprule
\textbf{Aspect} & \textbf{Verdict} & \textbf{Note} \\
\midrule
Password hashing & Correct & bcrypt, salted, never stores plain text. \\
Card handling & Correct & Redirect to hosted Checkout; cards never touch this app. \\
Server-side pricing & Correct & Client sends a slug, server owns the amount. \\
Secrets in environment & Correct & Though the fallback key is a footgun and \file{.env} is never read. \\
Price display integrity & Correct & Single source; verified across all nine planets. \\
Template structure & Good & One shared base, one looped card. \\
Session handling & \textbf{Broken} & Login is a global flag; every visitor shares one state. \\
Payment confirmation & \textbf{Broken} & \code{/success} is unauthenticated and unverified; no webhook, no orders table. \\
Input validation & Weak & Client-side only; several inputs return 500. \\
Data model & Weak & No unique email, wrong column type for the hash. \\
Tests & Absent & No test suite at all. \\
\bottomrule
\end{tabular}
\end{table}

\begin{keyidea}{What this project is, stated fairly}
It is a competent beginner web application that gets the two genuinely dangerous things
right (it does not store plain-text passwords and it does not touch card numbers) and then
gets the state management wrong in a textbook way. The value it carries into an interview
is not "I built a shop." It is: "I can show you a global mutable flag standing in for a
session, explain precisely why it breaks the moment a second visitor arrives, prove it
with three cookie jars, and tell you the six-line fix." That is a better story than a
flawless CRUD app, provided the flaws are stated first and by the author. Every defect in
this document is fixable in an afternoon, and the fact that they are localised is itself
evidence the structure underneath is sound.
\end{keyidea}

\end{document}
